% 2_objetivos.tex
\chapter{OBJETIVOS}

\section{OBJETIVO GERAL}
O presente trabalho tem como objetivo empregar cálculos baseados na DFT para investigar os aspectos químicos e estruturais dos tungstatos metálicos (ZnWO$_4$, NiWO$_4$ e CoWO$_4$). A análise buscará compreender a ativação das moléculas e O$_2$ e H$_2$O, visando elucidar os mecanismos de formação das ROS.
  
\section{OBJETIVOS ESPECÍFICOS}
\begin{itemize}
  \item Otimizar os parâmetros de cálculo de DFT para os tungstatos estudados;
  \item Criar modelos computacionais baseados na termodinâmica ab initio das principais terminações das superfícies expostas dos materiais estudados e determinar as superfícies mais estáveis, juntamente de suas energias de clivagem;
  \item Avaliar a estabilidade dinâmica dos modelos selecionados a partir da termodinâmica ab initio usando simulações de dinâmica molecular quântica, escolhendo os modelos mais promissórios;
  \item Estudar as propriedades eletrônicas das terminações por meio de cálculos de estrutura de bandas, DOS, função trabalho, carga Bader e diferença de densidade de carga;
  \item Analisar a interação de moléculas de O$_2$ e H$_2$O com as superfícies definidas, visando a ativação dessas moléculas para formação das ROS, elucidando se existem efeitos sinérgicos;
  \item Determinar os estados de transição que dominam no mecanismo de formação de ROS a partir da ativação das moléculas de O$_2$ e H$_2$O.
\end{itemize}



